\documentclass[11pt,a4paper]{article}
\usepackage[UTF8]{ctex}
\usepackage{amsmath,amssymb,bm,dsfont,booktabs,longtable,array}
\usepackage{geometry,fancyhdr,enumitem,needspace,xcolor,listings,graphicx}
\usepackage{xurl,hyperref}
\geometry{left=2.1cm,right=2.1cm,top=2cm,bottom=2cm}
\setlength{\headheight}{14pt}\setlength{\parindent}{2em}\setlength{\parskip}{0.32em}
\setlength{\emergencystretch}{2em}
\definecolor{courseblue}{HTML}{2457A6}
\hypersetup{colorlinks=true,urlcolor=courseblue,linkcolor=courseblue}
\setlist{leftmargin=2em,itemsep=0.3em}
\lstset{basicstyle=\ttfamily\footnotesize,columns=fullflexible,keepspaces=true,breaklines=true,showstringspaces=false,frame=single,backgroundcolor=\color{black!3},rulecolor=\color{black!15},aboveskip=1em,belowskip=.8em,xleftmargin=.3em,xrightmargin=.3em}
\newcommand{\code}[1]{\texttt{\small\detokenize{#1}}}
\newcommand{\R}{\mathbb R}
\newcommand{\relu}{\operatorname{ReLU}}
\newcommand{\task}[1]{\par\Needspace{5\baselineskip}\noindent\textbf{#1}\par}
\pagestyle{fancy}\fancyhf{}\lhead{AI3002 人工智能与机器学习基础}\rhead{LAB2：频率学习与输入表示}\cfoot{\thepage}
\title{\vspace{-1cm}\textbf{LAB2：神经网络如何学习\\不同频率的信号}}
\author{中国科学技术大学}\date{2026 年秋季学期\quad 学生版\\截止时间：2026 年 11 月 23 日 06:00}
\begin{document}\maketitle\thispagestyle{fancy}

个人完成，满分 100 分。使用 NumPy 手写梯度进行核验，使用 PyTorch 完成主实验；CPU 即可，无需 GPU。

\section*{提交与评分}
提交两项：报告\code{学号-姓名-LAB2.pdf}上传至 Blackboard（\href{https://bb.ustc.edu.cn}{bb.ustc.edu.cn}），单个\code{submission.py}上传至 OJ（\href{https://oj.temaurinum.moe}{oj.temaurinum.moe}）。不接收其他文件。报告建议 6--8 页；图表与运行记录直接放入报告。允许 AI 辅助，但须说明使用方式并给出独立核验，所有提交内容须能本人解释。

\begin{center}\begin{tabular}{p{3.3cm}rp{9cm}}\toprule
类别&分值&内容\\\midrule
OJ 代码&40&Fourier 映射 8；NumPy 前向 6；损失与反向 14；PyTorch 网络与参数对应 6；PyTorch 目标 6\\
OJ 最终性能&10&未公开信号实例上的统一重训与预测\\
报告&50&实现与三方核验 8；频率分析 12；联合选参 14；反例 10；二维迁移 4；复现与独立核验 2\\
环境创建&0&独立环境与检查，不计分\\\bottomrule
\end{tabular}\end{center}

报告由下述六项任务组成，各项分值合计为 $8+12+14+10+4+2=50$ 分。

\section{环境与文件（不计分）}
新建\code{ai3002-lab2}环境，不复用 LAB0/LAB1 环境。若尚未安装 Conda，请先按 LAB0 安装 Miniforge。

Windows 打开 PowerShell，进入实际保存的 LAB2 根目录；若仓库位于 LAB0 的示例位置，可执行第一行，否则自行替换路径。
\begin{lstlisting}
cd $HOME\ALML_public\LABs\LAB2
conda env create -f environment.yml
conda activate ai3002-lab2
python src/verify_env.py
\end{lstlisting}
\Needspace{9\baselineskip}macOS 打开 Terminal，进入实际目录后执行：
\begin{lstlisting}
cd ~/ALML_public/LABs/LAB2
conda env create -f environment.yml
conda activate ai3002-lab2
python src/verify_env.py
\end{lstlisting}
若 Shell 尚不能激活环境，Windows 先在 Miniforge Prompt 执行\code{conda init powershell}，macOS 执行\code{conda init zsh}，然后关闭并重开终端。在 VS Code 中运行 ``Python: Select Interpreter''，选择本实验环境。以后每次运行前只需激活，不必重新创建。环境已有同名环境时先检查其来源，不要直接删除个人环境。

检查成功显示\code{LAB2 environment ready}，并保存\code{results/environment.json}。确认其中解释器路径属于本实验环境。所有下述命令均从 LAB2 根目录运行；所有代码位于\code{src/}，数据位于\code{data/}，输出位于\code{results/}。只修改\code{src/submission.py}。代码、模型、运行日志和数据均不作为报告以外的额外附件提交。

\Needspace{8\baselineskip}\section{数据、网络与统一训练规则}
\task{2.1 数据及用途}
三个公开实例各有 256 个训练点和 256 个验证点，输入位于 $[-1,1)$。每个信号的形式为
\[
f(x)=\sum_{r=1}^{4}a_r\sin(\pi k_rx+\theta_r),\qquad
(a_1,a_2,a_3,a_4)=\frac{(1,0.5,0.3,0.2)}{\sqrt{(1^2+0.5^2+0.3^2+0.2^2)/2}}.
\]
每个 $k_r$ 分别从以下四组中抽取，$\theta_r$ 为随机相位：
\[
\{1,2\},\quad\{3,4,5\},\quad\{7,8,9,10\},\quad\{13,14,15,16,17,18\}.
\]
数据中的\code{freqs}记录该实例的四个频率。\code{X_train,y_train}用于训练，\code{X_val,y_val}用于选参。\code{X_grid,y_grid}是与二者不重合的 2048 点等距诊断网格，只用于画图、分析残差。诊断网格不是独立最终测试集，不得用其误差排名选参。不能在实现中直接编写目标函数、读取数据文件或猜测未公开相位。

另外提供 32 点稀疏训练数据，其标签加有标准差为 $0.15$ 的独立高斯噪声，以及一张 $64\times64$ 的课程生成灰度图。图像包含平滑区域、边缘与条纹，固定选出 1024 个观测像素；其余像素仅用于最后一次迁移评估。

\task{2.2 输入表示和网络}
对一维输入，Fourier 特征按以下顺序排列：
\[
\phi_s(x)=\bigl[x,\sin(\pi s2^0x),\ldots,\sin(\pi s2^5x),
\cos(\pi s2^0x),\ldots,\cos(\pi s2^5x)\bigr].
\]
二维时对每个坐标分别生成 13 维向量，再按坐标顺序拼接为 26 维。$s>0$ 为频率尺度，不是需要求梯度的模型参数。

设批量输入特征为 $F\in\R^{n\times d}$，标签为 $y\in\R^{n\times1}$。网络为
\[
\begin{aligned}
A_1&=FW_1+\mathds{1}_nb_1^\top,&H_1&=\relu(A_1),\\
A_2&=H_1W_2+\mathds{1}_nb_2^\top,&H_2&=\relu(A_2),\\
\widehat y&=H_2W_3+\mathds{1}_nb_3^\top.
\end{aligned}
\]
这里 $n$ 为批量大小，$d$ 为特征维度，$\mathds{1}_n$ 为全 $1$ 列向量；$\relu(t)=\max(0,t)$ 逐元素作用，零点导数在本实验定义为 $0$。主实验两个隐藏层均为 64 维，因此
\[
W_1\in\R^{d\times64},\quad W_2\in\R^{64\times64},\quad W_3\in\R^{64\times1},
\quad b_1,b_2\in\R^{64},\quad b_3\in\R^1.
\]
目标函数使用均方误差，注意没有额外的二分之一：
\[
J=\frac1n\sum_{i=1}^n(\widehat y_i-y_i)^2+
\frac\lambda2\sum_{\ell=1}^3\|W_\ell\|_F^2.
\]
其中 $\|W\|_F^2$ 表示全部矩阵元素的平方和，$\lambda\geq0$；偏置不受惩罚。

\task{2.3 训练协议（不得改变）}
统一使用全批次 Adam，学习率为 $\eta$，其余参数为
\[
\beta_1=0.9,\quad\beta_2=0.999,\quad\epsilon=10^{-8},\quad
\text{weight decay}=0.
\]
每次从课程指定初始化出发更新 800 次，采用最后一次参数；不提前停止，不热启动，不使用 AdamW。随机种子为 20261021；前两层权重方差为输入维数倒数的两倍，输出层为输入维数的倒数，偏置为零。课程提供初始化数组；PyTorch 必须原样装入。

主训练使用单线程 CPU 和单精度；梯度核验使用双精度。另提供宽度为 70 的原始坐标模型作为近似匹配的对照。

\Needspace{8\baselineskip}\section{任务一：实现与三方核验（报告 8 分）}
补全\code{submission.py}的五个计算接口，并在最终选参后填写配置接口。NumPy 部分不得调用 PyTorch 或自动微分；PyTorch 部分允许自动微分。预测及标签统一为二维列数组，不能混用\code{(n,)}与\code{(n,1)}导致广播。

\begin{longtable}{>{\raggedright\arraybackslash}p{4.7cm}>{\raggedright\arraybackslash}p{11cm}}\toprule
接口&要求\\\midrule
\code{fourier_features(x, scale)}&支持 NumPy 数组和 PyTorch Tensor，返回同后端、同精度、同设备的特征。Tensor 版本保持对输入的自动微分。\\
\code{numpy_forward(x, params)}&返回预测及自定义缓存。参数键为\code{W1,b1,W2,b2,W3,b3}；权重按“输入维数乘输出维数”存储。\\
\code{numpy_loss_backward(...)}&接收预测、标签、缓存、参数和正则化系数，返回标量目标及全部六个参数梯度。不得修改输入。\\
\code{build_torch_model(params)}&返回\code{Sequential(Linear,ReLU,Linear,ReLU,Linear)}，按给定参数尺寸创建并原样载入。\code{Linear.weight}需要转置。\\
\code{torch_objective(...)}&返回与 NumPy 相同且可微的标量目标。\\
\code{get_config()}&返回固定的\code{scale,lr,lam}三个数。\\\bottomrule
\end{longtable}
\begin{lstlisting}
python src/run.py check
\end{lstlisting}
公开检查使用远离 ReLU 零点的小网络，在相同参数与输入下，对零和非零正则化分别进行以下两组比较。记当前参数数组的 NumPy 解析梯度、PyTorch 自动微分梯度、中心有限差分梯度分别为 $g_{\rm np}$、$g_{\rm pt}$、$g_{\rm num}$。

对单个参数坐标 $u$，保持其余参数不变，中心有限差分为
\[
g_{\rm num}(u)=\frac{J(u+\varepsilon)-J(u-\varepsilon)}{2\varepsilon},\qquad
\varepsilon=10^{-6}.
\]
\textbf{第一组：NumPy 解析梯度与中心有限差分。}每个参数数组的归一化误差须满足
\[
E_{\rm num,np}
=\frac{\max|g_{\rm num}-g_{\rm np}|}
{\max\{1,\max|g_{\rm num}|,\max|g_{\rm np}|\}}
<10^{-6}.
\]
这里的最大值均遍历当前参数数组的全部元素。

\textbf{第二组：PyTorch 自动微分梯度与 NumPy 解析梯度。}脚本使用\code{np.allclose}逐元素检查，要求对每个元素索引 $i$ 都有
\[
|g_{{\rm pt},i}-g_{{\rm np},i}|
\leq 10^{-10}+10^{-9}|g_{{\rm np},i}|.
\]
其中绝对容差为 $10^{-10}$，相对容差为 $10^{-9}$；权重梯度先转置为 NumPy 的布局再比较。

\textbf{输出说明：}\code{check.json}中的\code{relative_errors}保存第一组的归一化误差。第二组通过断言检查，不单独保存误差数值；检查成功时\code{passed=true}表示两组均通过，第二组失败时会报\code{torch mismatch}。

\textbf{报告要求：}列出两种正则化下六组参数的第一组归一化误差，并记录第二组是否通过；说明为何检查点应避开 ReLU 零点。若修复过错误，给出一项错误、修正和核验记录；若首次通过，给出对批量平均或偏置惩罚的一项独立核验。

\section{任务二：观察频率学习过程（报告 12 分）}
\begin{lstlisting}
python src/run.py observe
\end{lstlisting}
脚本对三个公开实例运行原始坐标宽度 64、原始坐标宽度 70、Fourier 宽度 64 三个对照；固定 $s=1$、$\eta=10^{-3}$、$\lambda=0$。记录第 0、50、200、800 次更新的预测，并每 50 次更新记录训练、验证及诊断误差。

对等距诊断点 $x_j$、残差 $r_j=\widehat y(x_j)-f(x_j)$，工具计算
\[
\alpha_k=\frac2M\sum_{j=1}^M r_j\sin(\pi kx_j),\quad
\beta_k=\frac2M\sum_{j=1}^M r_j\cos(\pi kx_j),\quad
E_k=\frac{\alpha_k^2+\beta_k^2}{2}.
\]
这里 $M=2048$，$k$ 为该实例的一个目标频率。$E_k$ 衡量残差在该频率上的分量；只检查四个频率并不能覆盖全部误差，其和不必等于 MSE。

\textbf{报告要求：}选择一个实例展示三个对照的学习曲线，并汇总另两个实例是否支持相同结论；结合初始幅度，使用绝对误差或初值归一化误差说明哪些分量下降更快。讨论近似参数量匹配后结论是否改变。

\Needspace{8\baselineskip}\section{任务三：联合选参（报告 14 分）}
先运行六组基础对照：
\begin{lstlisting}
python src/run.py tune
\end{lstlisting}
尺度取 $\{0.25,1,4\}$，学习率取 $\{10^{-3},3\times10^{-3}\}$，正则化为 $0$。每个组合分别在三个实例训练，以三个验证 MSE 的平均值选出最佳组合。先保存并阅读这一轮结果，再在该最佳尺度、学习率下补选两个不同的正则化系数，范围为 $[10^{-6},10^{-1}]$。例如下面命令：
\begin{lstlisting}
python src/run.py tune --extra-lambdas 0.0001 0.01 --reason "Compare weak and strong penalties after reading coarse results"
\end{lstlisting}
将自己的观察、补选理由和实际命令写入报告；选值理由必须形成于查看补选结果之前。最终从六组基础与两组补选中，按平均验证 MSE 最小选择。相同时依次按尺度、学习率、正则化系数从小到大打破平局。脚本会输出\code{recommended_config}。补选阶段重新执行时保留最新两组，报告应如实说明曾进行的其他尝试。

\textbf{报告要求：}给出八行完整对照表，包括参数、训练是否成功、三个实例的验证误差和均值；说明为什么选这些补选值、最终如何选择、是否存在尺度过大却训练误差很低的情形。失败或未改善试验保留，不扣除其方法和分析部分的分数。不得根据密集网格、二维未观测像素或 OJ 结果反向调参。

\section{任务四：稀疏含噪反例（报告 10 分）}
\begin{lstlisting}
python src/run.py stress
\end{lstlisting}
在三个实例上使用提供的 32 点含噪训练数据。学习率固定为任务三选出的值，比较三种条件：中等尺度 $s=1$、高尺度 $s=4$、高尺度加正则化。前两者正则化为零。

第三种条件优先使用最终所选正则化。若最终配置为零正则化，则使用两个正补选值中平均验证误差较小者。不得根据反例结果再次选择正则化。

\textbf{报告要求：}汇总三种条件的含噪训练误差、干净验证误差和干净诊断误差；展示一个代表实例的预测图。回答：训练误差更小能否推出恢复干净函数更好？正则化是否在所有实例均有改善？哪些结论可由这些控制变量的对照支持，哪些仍不能？


\section{任务五：二维迁移（报告 4 分）}
\begin{lstlisting}
python src/run.py transfer
\end{lstlisting}
将最终所选尺度、学习率与正则化直接用于图像坐标回归，隐藏层宽度仍为 64、训练仍为 800 次。只在 1024 个观测像素上训练，不读取其他像素作为训练标签。此阶段只作一次迁移评估，不安排新的参数搜索。

\textbf{报告要求：}放入原图、观测掩码、重建图、绝对误差图，报告未观测像素 MSE；用简短文字指出误差主要出现于平滑区域、边缘还是条纹，并说明一维最优配置为何不保证二维最优。

\Needspace{8\baselineskip}\section{任务六：复现与独立核验（报告 2 分）}
将\code{results/tune.json}中的\code{recommended_config}原样填入\code{submission.py}的\code{get_config()}，然后执行：
\begin{lstlisting}
python src/run.py final
\end{lstlisting}
脚本核对配置与记录一致，重新训练并检查验证结果可复现。修改\code{get_config()}不会使实现核验失效；修改其他计算代码后须重新运行各阶段。数据校验失败应恢复课程数据，不应通过修改校验文件绕过。最终核验保存代码摘要、种子、配置、环境与耗时，供报告复核。

\begin{longtable}{>{\raggedright\arraybackslash}p{3cm}>{\raggedright\arraybackslash}p{6.2cm}>{\raggedright\arraybackslash}p{6.2cm}}\toprule
阶段&主要输出&报告用途\\\midrule
check&\code{check.json}&六组参数的三方梯度核验\\
observe&\code{observe.json}；每个实例的\code{history.json, predictions.png, frequencies.png}&频率学习过程与参数量对照\\
tune&\code{tune_coarse.json, tune.json, trials.csv}&粗比较、补选理由与最终八组结果\\
stress&\code{stress.json}及对应预测图&稀疏含噪反例\\
transfer&\code{transfer.json, transfer.png}&二维未观测像素误差与图像\\
final&\code{final.json}&最终配置、重训一致性、实现摘要\\\bottomrule
\end{longtable}
这些文件仅本地留存。\textbf{报告要求（共 2 分）：}
\begin{itemize}
\item \textbf{复现记录（1 分）：}在报告末尾记录软件版本、设备、完整运行命令、最终配置与总用时，使他人能按记录复现实验。
\item \textbf{独立核验记录（1 分）：}说明 AI 使用情况，并给出一项可复核的自检过程与结果；未使用 AI 也应给出自检。可引用任务一已有的独立核验记录，无需额外增加实验。
\end{itemize}
问题通过\href{https://github.com/Intelligent114/ALML_public/issues}{课程仓库 Issues}提出，注明 LAB2 和任务编号，不公开完整实现或个人信息。

\task{最终性能的简要规则}
OJ 使用八个未公开的一维信号实例，按统一 CPU 协议重新训练；你的配置在全部实例固定。参考基准使用相同架构与预算，固定 $s=0.25$、$\eta=3\times10^{-3}$、$\lambda=0$。汇总全部未见位置的 RMSE；不劣于基准得 10 分，否则按基准 RMSE 与提交 RMSE 的比值折算。训练失败或非有限输出时性能项为零，代码分独立评测。二维任务只计报告分。

只能提交通用计算实现，不能提交预训练权重、硬编码信号、读取外部数据或改变训练协议。最终性能仅占 10 分，报告按证据、方法和解释评分，未获得最好的单次曲线不妨碍获得分析分。

\end{document}
